% =====================================================================
% conclusiones.tex
% Responsabilidad única: responder las preguntas del laboratorio con las
% cifras de prueba, proponer acciones modestas, declarar limitaciones y
% trazar el trabajo futuro con la exportación real. Cifras: macros Fin*.
% =====================================================================
\section{Conclusiones, limitaciones y trabajo futuro}
\label{sec:conclusiones}

\subsection{Respuestas a las preguntas del laboratorio}

\begin{itemize}
    \item \textbf{Carga diaria:} una regresión lineal con calendario, puntos de
          muestreo programados y medias móviles predice la carga del día con
          un error absoluto medio de \FinMAEPrueba{} muestras sobre una media
          de \FinMediaPrueba{}, y una habilidad de \FinHabilidadPrueba{} frente
          a la media de los cinco días anteriores; aun así, el
          \FinDiasFueraPct{}\,\% de los días queda fuera de la tolerancia
          práctica.
    \item \textbf{Riesgo de no conformidad:} una regresión logística calibrada,
          con el umbral \FinUmbral{} que dicta la razón de costos, ordena las
          muestras con un ROC-AUC de \FinAUCPrueba{} y emite unas
          \FinAlertasSemana{} alertas por semana, de las cuales acierta una
          fracción de \FinPrecisionPrueba{} y que detectan una fracción de
          \FinRecallPrueba{} de las no conformes, con un ahorro del
          \FinAhorroPct{}\,\% frente a no alertar.
    \item \textbf{Falso positivo:} según los controles y el clasificador, uno de
          cada \FinUnoDeCadaFalso{} positivos sería falso, cifra que la
          fracción de reensayos conformes de prueba confirma
          (\FinReensayoConformeObservadoPct{}\,\% observada frente a
          \FinReensayoConformeEsperadoPct{}\,\% esperada), y los reensayos
          conformes se concentran en los recuentos que superan el límite en
          menos de \ClfReglaVecesLimite{} veces.
    \item \textbf{Grupos:} existen \FinAgrK{} perfiles de muestra y
          \FinAgrSegmentos{} segmentos de clientes, estables entre semillas
          pero con fronteras difusas (silueta de \FinAgrSilueta{}), útiles
          para describir, no para decidir de forma automática.
\end{itemize}

\subsection{Recomendaciones}

Las recomendaciones son deliberadamente modestas, porque las cifras vienen de
datos simulados y lo primero es comprobarlas con los reales:

\begin{itemize}
    \item \textbf{Planeación de turnos:} usar la predicción diaria como apoyo
          para asignar personal de recepción y siembra, con una revisión manual
          en las primeras semanas del año, donde el modelo sobreestima.
    \item \textbf{Piloto de alertas en recepción:} marcar las muestras que
          superan el umbral durante algunas semanas, sin cambiar el flujo,
          y medir con los datos reales la precisión, el recall y la carga de
          alertas por semana antes de adoptarlas.
    \item \textbf{Regla de reensayo:} priorizar el reensayo de los no conformes
          cuyo recuento quede por debajo de \ClfReglaVecesLimite{} veces el
          límite, que es donde se concentran los reensayos conformes.
    \item \textbf{Seguimiento de clientes:} revisar con el área comercial y de
          calidad el segmento con más no conformidad y más entregas tardías.
\end{itemize}

\subsection{Limitaciones}

\begin{itemize}
    \item \textbf{Datos simulados:} las métricas muestran que el procedimiento
          funciona de punta a punta y que no hay fugas, pero no miden la
          operación real; el generador fija relaciones que los modelos luego
          encuentran, así que con datos reales las habilidades pueden ser más
          modestas, y las razones de odds o los perfiles no deben leerse como
          hallazgos sobre el laboratorio.
    \item \textbf{Pocos casos en algunos ensayos:} \textit{Salmonella} y
          esterilidad tienen pocos positivos, de forma que su VPP y su
          fracción de reensayos conformes tienen intervalos anchos y, en
          prueba, no se interpretan.
    \item \textbf{Reensayo como referencia:} un reensayo conforme no equivale a
          un falso positivo, porque el reensayo también falla; con datos
          reales no existe otra verdad contra la cual comparar.
    \item \textbf{Siluetas bajas:} los grupos tienen fronteras difusas y la
          silueta del cribado no era comparable con la del estudio a fondo,
          por un cambio de preprocesamiento que el cribado debió compartir.
    \item \textbf{Intervalos bootstrap:} remuestrean filas como si fueran
          independientes, lo que en la serie diaria ignora la autocorrelación
          y hace que los intervalos de la regresión sean probablemente
          demasiado estrechos.
    \item \textbf{Un solo corte de prueba:} la prueba cubre un único periodo
          (\FinPruebaDesde{} a \FinPruebaHasta{}), así que las diferencias
          entre validación y prueba mezclan el desempeño del modelo con los
          cambios de nivel de la serie.
\end{itemize}

\subsection{Trabajo futuro con la exportación real}

El siguiente paso es conectar la exportación real, que solo requiere copiar
los archivos, completar el mapeo de columnas y reconstruir los notebooks,
revisar primero la tabla de cobertura y la auditoría de limpieza, y ajustar
las fechas de corte para que caigan dentro del periodo exportado. Después,
conviene repetir el cribado con el mismo preprocesamiento del estudio a
fondo, usar intervalos por bloques (bootstrap de bloques móviles) en la serie
diaria y, si el piloto de alertas confirma su utilidad, recalibrar el modelo
periódicamente, ya que la tasa de no conformidad cambia con la mezcla de
clientes.
